% Discussion: why verified self-training does not compound (distilled from GPT-6-Astra analysis, 2026-09-18).
% Static, hand-checked. \input by figures.tex.
\section{The discovery-to-assimilation bottleneck}

\paragraph{The mechanism: a discovery$\to$assimilation bottleneck (sharpening, not exploration).}
Our two strongest results jointly localise the failure. Best-of-$N$ search beating lineage training at matched
compute (lineage$-$bestof$N=-0.111\,[-0.138,-0.083]$ over 57 independent trajectories), together with the
large pass@$K\!\gg\!$pass@1 gap, shows
that \emph{useful solutions are already present in the sampling tail} but are not reliably produced in one
sample. The limiting pipeline is therefore
\[
\text{rare verified sample}\;\to\;\text{selected SFT example}\;\to\;\text{weight update}\;\to\;\text{reliably reusable gain},
\]
and it is \emph{lossy}: search can exploit a rare discovery directly, whereas training must compress it into a
shared parameterisation, preserve other skills, and have the gain survive later rounds. The lineage-vs-reset
null (inheriting updates does not beat restarting) is direct evidence \emph{against a useful cumulative learning
state}, not merely evidence that individual updates are small. This is the classic self-training signature:
selected-sample training concentrates probability on \emph{already-reachable} successes faster than it expands
what is reachable (cf.\ STaR, ReST, ReST$^{EM}$, expert iteration; and the RLVR capability-boundary debate,
Yue et al.\ 2025).

\paragraph{A formal fixed point --- with its limits stated precisely.}
Under an \emph{idealised} support-preserving rejection-conditioning operator $R$ (unlimited accepted samples,
exact per-task fitting, no cross-task transfer), exact solution-support coverage
$C_0(\pi)=\Pr_x[p_\pi(x)>0]$ is invariant and $R$ is idempotent after one step: $C_0(R\pi)=C_0(\pi)$ and
$R^2\pi=R\pi$. On tasks with positive success probability $R$ places all mass on accepted outputs; on $p{=}0$
tasks no accepted distribution exists. So an idealised verified-SFT loop is a \emph{one-step fixed point in
exact coverage}. \textbf{We are careful about what this does \emph{not} say.} (i) Neural SFT is not this
operator --- shared parameters can generalise and raise the probability of unsampled outputs, so real SFT need
not preserve support. (ii) \emph{Finite-budget} coverage is \emph{not} invariant:
$C_K(\pi)=\mathbb{E}_x[1-(1-p_\pi(x))^K]$ rises sharply under sharpening even when $C_0$ is fixed --- ``pass@$K$
coverage is a fixed point'' is false. (iii) Zero observed successes only bound $p\lesssim 3/K$; our
``unreachable mass'' is more precisely \emph{undiscovered mass at the evaluated budget}. (iv) A training-operator
fixed point is \emph{not} a parameter-space local optimum. We therefore report an \textbf{operational
self-training plateau / protocol-relative fixed point}, never an unqualified ``local optimum.''

\paragraph{Why verified search escapes where training does not.}
Search does not escape the policy's support; it escapes the \emph{one-sample} bottleneck and avoids the
assimilation bottleneck. For success probability $p$, $P(\text{hit in }N)=1-(1-p)^N$ makes extra draws
valuable when $p$ is small, and the verifier keeps the winner. Search (a) uses a discovery directly rather than
compressing it into weights, (b) allocates candidates per-task rather than via one shared update, and (c) does
not permanently contract the proposal policy. The precise statement: \emph{verified search accesses useful
existing tail capability more efficiently than the tested training loop consolidates it}. (This does not settle
lifetime economics --- a training update could amortise over enough future queries; that needs a separate
deployment-horizon study.)

\paragraph{Diversity contraction as a reinforcing lock (associational).}
Selected-sample SFT can drive a loop: current common winners $\to$ accepted set $\to$ higher probability on
those winners $\to$ fewer alternatives in the next accepted set. The loss is \emph{option value} --- rare
decompositions and kernel families that would enable \emph{later} discoveries. If undiscovered useful mass
$U_t$ shrinks, later rounds find fewer novel successes ($1-(1-U_t)^N\le NU_t$) even as pass@1 rises on familiar
ones. Higher retention among compounders and diversity collapse at large scale are \emph{consistent} with this
lock, but we flag them as association, not proof: low diversity can also reflect successful specialisation, and
the cleanest test is a matched-reward, matched-example, diversity-\emph{preserving} intervention.

\paragraph{The inverted-U, mechanistically.} Partition tasks into \emph{undiscovered} ($p{\approx}0$ at budget),
\emph{learnable frontier} ($0<p<0.2$), and \emph{near-ceiling} ($p$ high). Most self-training value comes from
the middle band. Small models: formation failures ($P(A\mid x)=P(F\mid x)P(A\mid F,x)$ with tiny $P(F)$) starve
the accepted set; mid-scale: the frontier is widest, so discoveries become signal (RSI peaks); large scale:
high coverage but little remaining headroom, so drift redistributes mass within already-covered behaviour
(``churn without gain'' --- supported descriptively; ``only superficial late-layer edits'' is \emph{not} yet
established, as parameter distance is not calibrated functional change).

\paragraph{What would break the plateau, and what our data predict.}
Two distinct breakouts are needed: \emph{accumulation} (retain/express more already-discoverable successes ---
could yield lineage$>$reset) and \emph{frontier expansion} (make new strategies operationally reachable ---
required once the frontier is exhausted). A productive loop needs
$\text{new discoveries}+\text{assimilation}+\text{retention}+\text{headroom}$ simultaneously. The single most
informative next experiment is a \textbf{coverage-injection $\times$ diversity-preserving} factorial: inject
stronger-model / repair / search-generated solutions on \emph{undiscovered} tasks and test whether inherited
gains then emerge --- separating ``lack of discoveries'' from ``loss of discoveries after learning.''
Exploration bonuses, novelty-selected accepted sets, prerequisite curricula, off-policy replay, and
search-amplified expert iteration are ranked candidates; each helps only the specific deficit it targets (e.g.\
entropy cannot fix formation; off-policy replay of redundant winners adds no frontier information).

\paragraph{Scope --- claims we do \emph{not} make.} We do not claim the models' reachable set is mathematically
invariant, that the weights sit at a genuine local optimum, or that RL/RLVR cannot expand capability. We claim,
with \emph{moderate} evidence for the null and a significant positive respectively: under verified
rejection-sampling SFT at the tested scales/budget, carrying the lineage forward does not beat a matched reset
(trajectory-level TOST equivalence at $\delta{=}0.05$, $\mathrm{BF}_{01}{=}5.6$ over 57 independent
trajectories); verified search is a better use of the same compute (CI excludes zero); and the binding internal signals are coverage/formation (small), a mid-scale learnable
frontier, and drift-without-gain plus diversity contraction (large). The synthesis: \emph{expert iteration
without sufficient expert amplification, plus incomplete retention of the apprentice's useful tail.}
